\section{产品评测清单：从 Wow 到 Retention}

前几章讨论了妙鸭、AI Native 和 One Way Door，但产品判断还需要更可操作的清单。张月光在访谈里多次把“好看”“惊艳”和“长期正确”分开：一个产品可以很 wow，却不是单向门；可以短期爆发，却不能长期留存；可以技术很强，却没有用户心智。

\subsection{四层评测}

\begin{longtable}{p{0.20\textwidth}p{0.36\textwidth}p{0.35\textwidth}}
\toprule
层次 & 要问的问题 & 妙鸭 / Dokie 的启发 \\
\midrule
技术成立 & 这代模型是否能稳定做出结果？ & 妙鸭选择当时相对成熟的图像技术；Dokie 要验证 Agent 是否能真实完成创作表达。 \\
效果函数 & 用户为什么觉得它好？ & 妙鸭是“真、像、美”；Dokie 是“突破能力边界”。 \\
单向门 & 用户用了后是否回不去旧方法？ & 妙鸭是强体验但不一定永久替代；Agent 若能稳定干活，更可能成为单向门。 \\
心智与留存 & 用户会用什么词记住它？是否足够高频？ & PPT 太窄，创作与表达更宽；AI 朋友则需要长期关系感。 \\
\bottomrule
\end{longtable}

\begin{importantbox}{从爆款到长期产品}
爆款通常证明技术、传播和需求击中过一次；长期产品还要证明心智、频次、留存、复购和能力进化。AI 产品尤其要警惕“只证明了 wow，没有证明单向门”。
\end{importantbox}

\subsection{为什么“我自己会用”仍然重要}

张月光说，他很难做自己完全不会用的产品判断。这个标准并不适用于所有产品经理，但在 AI 早期尤其重要，因为很多指标还没稳定，用户也说不清楚自己要什么。创始人自己是否真的离不开这个新解法，常常是最早的定性信号。

\begin{warningbox}{自用不是万能验证}
自用能捕捉早期直觉，但不能替代用户分层、真实场景和商业验证。它适合作为启动信号，不适合作为最终证明。
\end{warningbox}

\subsection{本章小结}

AI 产品评测要同时看技术、效果函数、单向门和用户心智。妙鸭证明张月光能做爆款，Dokie 要证明他能做长期正确的 Agent。

